% =============================================================================
% 008 / 068
% Status: raw
% =============================================================================
% Notes:
%
% Source question:
% 上海市區話已經無法搶救，是否能從 松江 嘉定，浦東崇明等地有殘留原生詞彙語法的重新構建出一套出來？這些郊縣，那個語法標記詞彙儲存最完整？
% =============================================================================

\chapter[上海市區話已經無法搶救，是否能從 松江 嘉定，浦東崇明等地有殘留原生詞彙語法的重新構建出一套出來？這些郊縣，那個語法標記詞彙儲存最完整？]{上海市區話已經無法搶救，是否能從 松江 嘉定，浦東崇明等地有殘留原生詞彙語法的重新構建出一套出來？這些郊縣，那個語法標記詞彙儲存最完整？}
\qaanswer{
上海市區話本身在形成過程中就一種藍青官話。因為本身就沒有標準，語法體系又都是官話一脈的基本上沒辦法拯救，放到50年前也救不了。上海郊縣語法標記保留的都還可以，特別是松江，有很多熱心的同胞在整理和紀錄語法、詞彙，保留了大量的文獻記載。
}

